% Part: second-order-logic
% Chapter: syntax-and-semantics
% Section: terms-formulas

\documentclass[../../../include/open-logic-section]{subfiles}

\begin{document}

\olfileid[pt]{sol}{syn}{frm}

\olsection{Termos e \printtoken{P}{formula}}

Assim como na lógica de primeira ordem, as expressões da lógica de segunda
ordem são construídas a partir de um vocabulário básico que contém
\emph{!!{variable}s}, \emph{!!{constant}s}, \emph{!!{predicate}s} e, às
vezes, \emph{!!{function}s}. A partir deles, junto com conectivos lógicos,
quantificadores e sinais de pontuação como parênteses e vírgulas, formam-se
\emph{termos} e \emph{!!{formula}s}. A diferença é que, além de variáveis
para objetos, a lógica de segunda ordem também contém variáveis para
relações e funções, e permite quantificar sobre elas. Portanto, os símbolos
lógicos da lógica de segunda ordem são os da lógica de primeira ordem,
acrescidos dos seguintes:

\begin{enumerate}
\item Um conjunto !!{denumerable}s de !!{variable}s de relação de segunda
  ordem de cada aridade~$n$: $\Obj V_0^n$, $\Obj V_1^n$, $\Obj V_2^n$,
  \dots
\item Um conjunto !!{denumerable}s de !!{variable}s de função de segunda
  ordem: $\Obj u_0^n$, $\Obj u_1^n$, $\Obj u_2^n$, \dots
\end{enumerate}

Assim como usamos $x$, $y$, $z$ como metavariáveis para as variáveis de
primeira ordem $\Obj v_i$, usaremos $X$, $Y$, $Z$ etc. como metavariáveis
para $\Obj V_i^n$, e $u$, $v$ etc. como metavariáveis para~$\Obj u_i^n$.

\begin{explain}
Os símbolos não lógicos de uma linguagem de segunda ordem são especificados
da mesma maneira que os de uma linguagem de primeira ordem: enumerando suas
!!{constant}s, !!{function}s e !!{predicate}s.

Na lógica de primeira ordem, a !!{identity}~$\eq$ costuma estar incluída.
Na lógica de primeira ordem, os símbolos não lógicos de uma linguagem
~$\Lang{L}$ são fundamentais para que possamos expressar algo interessante.
É claro que há !!{sentence}s que não usam símbolos não lógicos, mas apenas
com~$\eq$ é difícil dizer algo interessante. Na lógica de segunda ordem,
como temos um suprimento ilimitado de variáveis de relação e de função,
podemos dizer tudo o que podemos dizer em uma linguagem de primeira ordem
mesmo sem um suprimento especial de símbolos não lógicos.
\end{explain}

\begin{defn}[Termos de segunda ordem]
O conjunto dos \emph{termos de segunda ordem} de~$\Lang L$, $\TrmSOL[L]$,
é definido acrescentando a \olref[fol][syn][frm]{defn:terms} a cláusula
\begin{enumerate}
\item Se $u$ é uma variável de função de $n$ argumentos e $t_1$, \dots,
  $t_n$ são termos, então $\Atom{u}{t_1, \ldots, t_n}$ é um termo.
\end{enumerate}
\end{defn}

\begin{explain}
Assim, um termo de segunda ordem se parece exatamente com um termo de
primeira ordem, exceto pelo fato de que, onde um termo de primeira ordem
contém !!a{function}~$\Obj{f^n_i}$, um termo de segunda ordem pode conter em
seu lugar uma variável de função~$\Obj{u^n_i}$.
\end{explain}

\begin{defn}[\usetoken{s}{formula} de segunda ordem]
O conjunto das \emph{!!{formula}s de segunda ordem}~$\FrmSOL[L]$ da
linguagem~$\Lang L$ é definido acrescentando a
\olref[fol][syn][frm]{defn:formulas} as cláusulas
\begin{enumerate}
\item Se $X$ é uma variável de predicado de $n$ argumentos e $t_1$, \dots,
  $t_n$ são termos de segunda ordem de~$\Lang L$, então
  $\Atom{X}{t_1,\ldots, t_n}$ é uma !!{formula} atômica.

\tagitem{prvAll}{Se $!A$ é !!a{formula} e $u$ é uma variável de função,
  então $\lforall[u][!A]$ é !!a{formula}.}{}

\tagitem{prvAll}{Se $!A$ é !!a{formula} e $X$ é uma variável de predicado,
  então $\lforall[X][!A]$ é !!a{formula}.}{}

\tagitem{prvEx}{Se $!A$ é !!a{formula} e $u$ é uma variável de função,
  então $\lexists[u][!A]$ é !!a{formula}.}{}

\tagitem{prvEx}{Se $!A$ é !!a{formula} e $X$ é uma variável de predicado,
  então $\lexists[X][!A]$ é !!a{formula}.}{}
\end{enumerate}
\end{defn}

\end{document}

